\label{sec:discussion}

\subsection{What the experiment sequence taught}
\label{sec:discussion-arc}

The motivating idea, use onboard vision to shutter only on clear targets, is simple to state but hard to train.
The fourteen experiments were not a single attempt to beat the baseline; they were a structured search for reward and action-space configurations that make learning possible.
Three lessons emerged.
Per-experiment returns and shutter counts are reported in Sections~\ref{sec:experiments} and~\ref{sec:results}; this subsection interprets why the levers mattered.

\paragraph{1.\ Reward structure before algorithm.}
Neither threshold tuning, encoder architecture, algorithm choice, nor network width was sufficient on its own.
Learning emerged only after the reward told the agent what mattered at the action clock: sparse capture credit, a budget-exhausted penalty that stops shutter spam, and removal of a redundant within-budget waste term.
Reward engineering was the dominant design lever; swapping algorithms without that structure did not unlock selective shuttering.

\paragraph{2.\ Action-space design before reward fine-tuning.}
Delegating rate control to the onboard pointing controller (vector mode) is structurally important relative to raw torque commands.
Vector mode removes safe-mode lockouts in the experiment runs (Section~\ref{sec:results-safemode}), resolving a credit-assignment problem that no reward penalty could fix cleanly.
This is a design result: the agent should command a pointing offset, not raw wheel torque, for cloud-aware scheduling in this environment (Section~\ref{sec:discussion-architecture}).

\paragraph{3.\ Algorithm stability as a prerequisite.}
An unconstrained MPO trust-region update blocked the MPO branch until the dual-variable bug was corrected.
After that fix, vector-mode MPO learned without a further reward redesign.
Implementation correctness can masquerade as algorithm failure; that check comes before comparing reward or action-space variants.

\subsection{Answer to the research question}
\label{sec:discussion-answer}

Cloud-aware autonomous scheduling is not trivial to train in this simulation.
A minimal viable training stack requires at least:
\begin{itemize}
  \item a sparse capture-credit reward with a budget-exhausted shutter penalty;
  \item vector pointing mode (nadir-relative offset) rather than direct torque commands;
  \item a correctly implemented MPO trust-region update, or SAC as the more forgiving default;
  \item baseline-policy warmup to seed the replay buffer with rare positive-credit transitions.
\end{itemize}
With that stack, both SAC and MPO learn selective shuttering absent in the deterministic baseline.
Several configurations that looked reasonable (threshold tuning, network width, safe-mode penalties, torque effort in vector mode) failed to produce learning or made it worse.
Eval return on seed~7 exceeded the matched average baseline return for the best SAC and MPO configurations, but the evaluation protocol is too narrow to claim robust mission superiority; the durable output of the semester is the mapped design space and the list of rejected shortcuts.

\subsection{Control hierarchy and action interface}
\label{sec:discussion-architecture}

Development did not begin with the interface the experiments ultimately recommend.

\paragraph{From torque to safety arbitration.}
The first experimental interface exposed raw reaction-wheel torque to the learning agent.
Unconstrained policies could command aggressive slews; in simulation this produced sustained rotation that would be unacceptable on flight hardware.
Review with system engineers motivated treating attitude safety as non-negotiable: every agent command, whether torque or pointing request, passes through a safety controller that can taper, reject, or replace commands with a safe-mode recovery sequence (Section~\ref{sec:safety}).
That stack is shared by the deterministic baseline and all learned policies.

\paragraph{Action modes tested.}
Three action interfaces were explored:
\begin{itemize}
  \item \textbf{Torque mode:} the policy commands wheel torque directly (Exps~1--5, 8, 10).
  \item \textbf{Vector mode:} the policy commands a nadir-relative pointing offset; the OBC converts it to PD torque (Exps~4, 7, 9, 11--13).
  \item \textbf{Target selection:} a higher-level discrete choice of which target to engage (Exp~14).
    Started as a hypothesised next simplification after torque and vector; the run did not complete, so whether this interface is easier to train remains open.
\end{itemize}

\paragraph{Segregation of control layers.}
The experiment arc supports separating what the agent should learn from what the OBC should execute.
Learning precise torque profiles or continuous pointing offsets is not the primary scheduling problem in this environment: the agent must decide \emph{when} to expose the shutter given cloud feedback, while rate control and maneuver execution belong downstream.
Vector mode reflects that split and improved learning in practice (Exps~4, 11).
A torque-trained agent might in principle discover rich low-level strategies, for example multiple captures within one forward-motion-compensation maneuver.
It is a further \emph{hypothesis}, not a result of this report, that raising the action interface to explicit target selection (plus shutter timing) would simplify credit assignment further; Exp~14 was begun on that conviction but time ran out before it could be validated, and implementation or gradient issues in that branch have not been ruled out.
Energy-optimal wheel commands remain a distinct optimisation problem; they should be trained or tuned in a separate low-level loop rather than entangled with cloud-aware target scheduling.

\paragraph{Intended end state (working hypothesis).}
A plausible long-term stack is three levels: mission intent from the ground or onboard planner, a vision-driven scheduling policy (target and shutter decisions), and OBC execution with PD pointing and hard safety limits.
Vector mode is the highest-level interface validated here; target selection is the candidate next step named above.
This semester maps the scheduling layer under vector actions in a simplified 2D simulator; the lowest layer is modelled but not co-trained for efficiency.

\subsection{Rejected paths}
\label{sec:discussion-rejected}

Table~\ref{tab:rejected} summarises design choices that were tested and ruled out.

\begin{table}[htbp]
  \centering
  \caption{Rejected experimental directions with rationale.}
  \label{tab:rejected}
  \begin{tabular}{p{0.22\linewidth}p{0.34\linewidth}p{0.36\linewidth}}
    \toprule
    Experiment & Design choice tested & Why rejected \\
    \midrule
    --- & MPC pointing controller & Superseded by onboard PD pointing; lower scope and comparable results \\
    --- & Latent-only dense reward & Dense reward tested at $10\times$ applied scale; did not unblock MPO \\
    Exp~1 & Shutter threshold 0.5 $\to$ 0.9 & Policy saturates shutter output regardless of threshold \\
    --- & 15\,s capture credit window & Latent bursts visible; applied capture remains zero \\
    Exp~5 & MPO network width ablation & Identical collapse at all widths; capacity not the bottleneck \\
    Exp~9 & Within-budget waste penalty (on) & Redundant with apply-credit sparsity; penalising same failure mode twice \\
    Exp~10 & Safe-mode penalty & Correlation with aggressive pointing suppresses useful captures \\
    Exp~12 & Torque effort in vector mode & Pointing-controller torque $\neq$ policy action; conflicting gradient \\
    \bottomrule
  \end{tabular}
\end{table}

\subsection{Limitations}
\label{sec:limitations}

\paragraph{Two-dimensional simulation.}
The entire dynamics model, camera model, and reward signal operate in a single orbit plane.
Real three-axis attitude control involves cross-coupling between roll, pitch, and yaw;
disturbance torques from solar pressure, magnetic field, and gravity gradient; and
three-dimensional off-nadir imaging geometry.
The 2D model provides a tractable learning environment but the policies are not directly
transferable to hardware.

\paragraph{Single-pass, single-orbit mission.}
Each episode models one overflight of a fixed 50-target corridor.
Revisit scheduling, downlink windows, and multi-orbit memory are not represented.
The learned policy has no access to observations from prior passes.

\paragraph{Simulation-only evaluation.}
No hardware-in-the-loop or software-in-the-loop tests were conducted.
The camera model is a geometric ray tracer with an analytic smear formula; radiometric calibration,
MTF, and sensor noise are absent.
Generalisation from this environment to a real EO instrument is an open question.

\paragraph{Fixed-seed evaluation.}
All primary eval returns are computed on two episodes at seed~7.
The train-to-eval gap visible in several experiments (e.g.\ Exp~9 train $+169.9$ vs eval $+106.4$)
may partially reflect evaluation seed dependence.
A robust evaluation across $\ge$10 seeds is recommended before citing the best numbers as
reliable performance estimates.

\paragraph{Budget and episode length.}
The 10-capture budget over 516 steps is an engineered constraint that simplifies mission planning
but may not represent realistic EO tasking where capture rate, data volume, and ground station contact
are coupled objectives.

\subsection{Deferred extensions}
\label{sec:discussion-deferred}

The following directions are documented as promising but were not completed within the project:

\begin{itemize}
  \item \textbf{Modular encoder rerun (deferred Exp~6).}
    Once learning is established, the vector-compressor encoder variant from Exp~2 should be retested.
    The high-dimensional observation with structured bearing and mask fields is
    a candidate for entity-based architectures \citep{gorishniy2022revisiting}.

  \item \textbf{MPO with Exp~9 reward config.}
    Exp~11 reached $+45.5$ without the waste-penalty-off configuration used in Exp~9.
    A direct MPO run at Exp~9 reward settings would quantify the algorithm gap more cleanly.

  \item \textbf{Target-selection interface (Exp~14 rerun).}
    Hypothesis: discrete target choice plus shutter timing delegates pointing entirely to the OBC and may simplify learning relative to vector mode.
    Exp~14 began testing this idea but crashed before a conclusive result; a rerun must also rule out implementation or gradient bugs in that branch, not only extend training time.

  \item \textbf{Three-axis attitude extension.}
    Moving from 2D to 3D dynamics is the prerequisite for any hardware-facing transfer.
    The simulation architecture supports this extension via replacement of the single-axis
    RW model with a 3-wheel assembly and attitude quaternion propagation.

  \item \textbf{Real imagery and radiometric reward.}
    Replacing the geometric smear quality metric with a model-based image quality score
    derived from sensor noise and scene contrast would couple the reward more directly
    to downstream science value.
\end{itemize}
